\section{Software: MANTA}
\label{sec:WP7:MANTA:software}

\begin{table}[h!]
    \centering
    { \setlength{\parindent}{0pt}
    \def\arraystretch{1.25}
    \arrayrulecolor{numpexgray}
    {\fontsize{9}{11}\selectfont
    \begin{tabular}{!{\color{numpexgray}\vrule}p{.4\textwidth}!{\color{numpexgray}\vrule}p{.6\textwidth}!{\color{numpexgray}\vrule}}
        \rowcolor{numpexgray}{\rule{0pt}{2.5ex}\color{white}\bf Field} & {\rule{0pt}{2.5ex}\color{white}\bf Details} \\
        \rowcolor{white}\textbf{Consortium} & \begin{tabular}{l}
CEA + consortium in development (see EUROPLEXUS)\\
\end{tabular} \\
        \rowcolor{numpexlightergray}\textbf{Exa-MA Partners} & \begin{tabular}{l}
CEA\\
\end{tabular} \\
        \rowcolor{white}\textbf{Contact Emails} & \begin{tabular}{l}
olivier.jamond@cea.fr\\
\end{tabular} \\
        \rowcolor{numpexlightergray}\textbf{Supported Architectures} & \begin{tabular}{l}
CPU Only\\
\end{tabular} \\
        \rowcolor{white}\textbf{Repository} & \href{None}{None} \\
        \rowcolor{numpexlightergray}\textbf{License} & \begin{tabular}{l}
None\\
\end{tabular} \\
        \bottomrule
    \end{tabular}
    }}
    \caption{WP7: MANTA Information}
\end{table}

\subsection{Software Overview}
\label{sec:WP7:MANTA:summary}

MANTA is CEA's explicit dynamics framework for high-strain-rate and impact studies. The project is transitioning legacy
workflows into the Exa-MA methodology: WP7 supports the creation of container recipes, CI smoke tests, and benchmarking harnesses
for representative MANTA kernels. Because the production sources are not yet open, WP7 focuses on publishing derived
mini-applications that mirror the solver stack and data dependencies while remaining distributable to partners.

\begin{table}[h!]
    \centering
    { 
        \setlength{\parindent}{0pt}
        \def\arraystretch{1.25}
        \arrayrulecolor{numpexgray}
        {
            \fontsize{9}{11}\selectfont
            \begin{tabular}{!{\color{numpexgray}\vrule}p{.25\linewidth}!{\color{numpexgray}\vrule}p{.6885\linewidth}!{\color{numpexgray}\vrule}}
    
    \rowcolor{numpexgray}{\rule{0pt}{2.5ex}\color{white}\bf Features} &  {\rule{0pt}{2.5ex}\color{white}\bf Short Description }\\ 
    
\rowcolor{white}    Energy & Shock hydrodynamics and structural integrity assessments for energy infrastructures (reactor vessels, cryogenic tanks). \\
\rowcolor{numpexlightergray}    mini-apps & Linear elasticity and contact mini-apps extracted from MANTA kernels to validate WP7 packaging and CI pipelines. \\
\rowcolor{white}    proxy-apps & Coupled fluid-structure-impact proxies aggregating three-stage workflows (mesh generation, explicit dynamics, post-processing). \\
\end{tabular}
        }
    }
    \caption{WP7: MANTA Features}
\end{table}


\subsection{Parallel Capabilities}
\label{sec:WP7:MANTA:performances}


MANTA relies on distributed-memory parallelism (MPI) and task-based OpenMP regions within the explicit dynamics loop. Hybrid
MPI+OpenMP configurations are routinely executed on CEA's TGCC clusters, with current qualification runs reaching 4096 CPU cores.
GPU offloading pilots using OpenACC are underway for contact detection kernels, coordinated with WP3. The WP7 integration layers
MANTA with shared mesh tools (ParMmg) and I/O standards (HDF5/Exodus) delivered by WP1, ensuring compatibility across the
Exa-MA application portfolio.


\subsection{Initial Performance Metrics}
\label{sec:WP7:MANTA:metrics}

Because MANTA remains closed-source, WP7 publishes anonymised benchmark metadata derived from representative kernels. The first
dataset (pending DOI assignment) covers a tank-impact scenario with 30~million elements. Executions on TGCC's IRENE compute
partition deliver 0.41~s/step on 2048 cores (82\% weak-scaling efficiency from 256 cores). Energy traces indicate that NUMA-aware
thread pinning reduces power draw by 11\%. Key challenges include long preprocessing chains in legacy scripts and limited
observability for I/O throughput, which are being addressed by wrapping the workflow with WP7's ReFrame runners and Darshan
instrumentation.

\subsubsection{Benchmark \#1}
\begin{itemize}
    \item \textbf{Description:} Tank-impact explicit dynamics case on IRENE Rome nodes (up to 2048 cores) assessing weak scaling and energy footprint.
    \item \textbf{Benchmarking Tools Used:} ReFrame automation, Extrae MPI profiling, Darshan I/O tracing, and WP7 energy plug-in for RAPL measurements.
    \item \textbf{Input/Output Dataset Description:}
        \begin{itemize}
            \item \textbf{Input Data:} Geometry and material parameter pack (2.1~GB, Exodus format)---DOI registration in progress via CEA data portal.
            \item \textbf{Output Data:} Time-resolved stress/strain fields (HDF5) and telemetry logs exported by ReFrame for post-processing.
            \item \textbf{Data Repository:} Interim storage on TGCC object store with mirrored metadata in the Exa-MA knowledge base pending public release.
        \end{itemize}
    \item \textbf{Results Summary:} 82\% weak-scaling efficiency (256\,$\rightarrow$\,2048 cores); energy-to-solution improved by 11\% after NUMA pinning; peak bandwidth 45~GB/s during mesh updates.
    \item \textbf{Challenges Identified:} Legacy preprocessing pipeline is still serial; WP7 is prioritising parallel mesh generation and staged I/O to mitigate queue times.
\end{itemize}

\subsection{12-Month Roadmap}
\label{sec:WP7:MANTA:roadmap}

\begin{itemize}
    \item \textbf{Data improvements:} Finalise anonymisation and DOI publication for the tank-impact benchmark dataset, including documentation on input preparation (Q1~2026).
    \item \textbf{Methodology application:} Deliver a reproducible MANTA mini-app container with automated ReFrame smoke tests covering MPI+OpenMP and experimental OpenACC builds (Q2~2026).
    \item \textbf{Results retention:} Integrate MANTA metrics into the WP7 Grafana dashboards and archive quarterly benchmark summaries in the Exa-MA knowledge base (starting Q3~2026).
\end{itemize}
